\documentclass{jsarticle}
\usepackage{amsmath}
\usepackage{amssymb}
\begin{document}
\title{アカシックレコードでの薬の調合についての証明のレポート}
\author{Masaaki Yamaguchi}
\date{}
\maketitle


       ドリエルは紅茶を合わせると、ガンマ・アミノ酪酸になり、
  抑うつや、躁を緩和する。松果体にも作用して、鎮静させる。
  マカは、利尿作用があり、腎臓透析に有効である。
  カリウムは、排尿抑制効果があり、尿が出にくくなる。
  フェニトインは、自問自答する効果があり、アカシックレコードに直結する。
  リボトリールは、このアカシックレコードの副作用止であり、
  体内のリボソームの複製を止める。
  アメルは、磁性体発生で、体と精神のアカシックレコードの正常化に効果がある。
  エチゾラムとブロチゾラムは、瞑想法の極意である。

\end{document}
